\subsection{RING HOMOMORPHISMS}

DEFINITION 2. Let A and B be two rings. A morphism, or homomorphism, of A into B is any mapping f of A into B satisfying the relations:

\[
f (x + y) = f (x) + f (y), \quad f (x y) = f (x). f (y), \quad f (1) = 1,\tag{17}
\]

for all \(x, y\) in A.

The composition of two ring homomorphisms is a ring homomorphism. Let A and B be two rings and f a mapping of A into B; for f to be an isomorphism, it is necessary and sufficient that it be a bijective homomorphism; in that case, \(\overline{f}^{1}\) is a homomorphism of B into A. A homomorphism of a ring A into itself is called an endomorphism of A.

Let \(f: A \to B\) be a ring homomorphism. The mapping \(f\) is a homomorphism of the additive group of \(A\) into the additive group of \(B\) ; in particular, \(f(0) = 0\) and \(f(-x) = -f(x)\) for all \(x \in A\) . The image under \(f\) of an invertible element of \(A\) is an invertible element of \(B\) and \(f\) induces a homomorphism of the multiplicative group of \(A\) into the multiplicative group of \(B\) .

Examples. (1) Let A be a ring. It is immediately seen that the mapping \(n \mapsto n\) . 1 of Z into A is the unique homomorphism of Z into A. In particular, the identity mapping of Z is the unique endomorphism of the ring Z.

In particular, take A to be the endomorphism ring of the additive group Z (no. 3, Example VI). The mapping \(n \mapsto n\) . 1 of Z into A is an isomorphism of Z onto A by the very construction of multiplication in Z (§ 2, no. 6).

\begin{enumerate}
\def\labelenumi{(\arabic{enumi})}
\setcounter{enumi}{1}
\tightlist
\item
  Let \(a\) be an invertible element of a ring \(A\) . The mapping \(x \mapsto axa^{-1}\) is an endomorphism of \(A\) for
\end{enumerate}

\[
\begin{array}{r l} a (x + y) a ^ {- 1} & = a x a ^ {- 1} + a y a ^ {- 1}, \\ a (x y) a ^ {- 1} & = (a x a ^ {- 1}) (a y a ^ {- 1}). \end{array}
\]

It is bijective, for the relation \(x' = axa^{-1}\) is equivalent to \(x = a^{-1}x'a\) . It is therefore an automorphism of the ring A, called the inner automorphism associated with a.

